\documentclass[11pt]{article}
\usepackage[margin=0.75in]{geometry}
\usepackage{fontspec}
\usepackage{setspace}
\usepackage{hyperref}
\hypersetup{colorlinks=true, urlcolor=blue, linkcolor=black}
\pagestyle{empty}
\setlength{\parskip}{0.65em}
\setlength{\parindent}{0pt}
\setstretch{0.93}
\setmainfont{Times New Roman}

\begin{document}

\begin{flushleft}
Ke Cheng HSU (Clio) \\
Department of Economics, The Hong Kong University of Science and Technology \\
Clear Water Bay, Kowloon, Hong Kong $\cdot$ +852 6840 9753 \\
\href{mailto:khsu@connect.ust.hk}{khsu@connect.ust.hk} $\cdot$ \href{https://cliohsu1997.github.io/}{cliohsu1997.github.io}
\end{flushleft}

24 November 2026

\vspace{0.1em}
Dr Prateek Bansal \\
Department of Civil and Environmental Engineering \\
College of Design and Engineering \\
National University of Singapore \\
Kent Ridge, Singapore

\vspace{0.1em}
Dear Dr Bansal,

I am writing to apply for the Research Fellow (DATA FUSION) position in your group (NUS Careers job 34585). I am a Ph.D.\ candidate in Economics at the Hong Kong University of Science and Technology. I expect to graduate in June 2027 and can join as early as February 2027. My fields are Empirical Industrial Organization and Digital Economics. The posting asks for discrete choice modelling, optimization, and econometric programming to fuse survey data with passively collected data---skills that are central to my structural empirical work.

My job market paper, ``Designing Return Policy in E-Commerce: A Trade-off between Purchase and Return Efficiency,'' estimates a structural discrete-choice style model of purchase, return-shipping-insurance take-up, and returns using large Taobao transaction data. I recover demand and uncertainty parameters by GMM with bootstrap inference, then use counterfactual optimization over policy menus to rank return policies by seller profit, consumer surplus, and social welfare. The project is built as a reproducible Python / R / C++ pipeline on AWS and HPC, with Monte Carlo validation of the estimator---the same validate-on-controlled-data habit the posting describes for open-source datasets.

A second strand of my work is closer still to multi-source data fusion. With Kohei Kawaguchi and Kei Kawai, I construct title-by-market-by-month panels from Japanese bookstore point-of-sale deliveries, sales, and returns (tens of billions of records), merge them with scraped product attributes and local demographics, and estimate a dynamic oligopoly model of pricing and inventory under resale price maintenance. That project combines ``survey-like'' attribute information with high-frequency passively collected retail records, which is the methodological analogue of fusing survey responses with smart-card or other sensor streams. With Pak Hung Au, I also have a theory paper, ``Product Design with Heterogeneous Search Costs,'' on how firms choose product design when consumers differ in search costs. With Ce Bian and Liang Zhao, I study pass-through when products can be returned, and show that a tax can lower the equilibrium price and raise welfare by reducing the return rate.

I am proficient in programming econometric models in Python and R (and in C++ for equilibrium simulation), and I routinely solve optimization problems inside estimation and counterfactual routines. I write and present clearly: I was a teaching assistant from 2020 to 2025 (average evaluation 4.72/5) and designed and lectured the MSc mathematics camp in 2026 (4.53/5). I am prepared to pursue independent and collaborative research in your group, to help write grant proposals, and to support supervision of doctoral and postgraduate students on fixed-term appointment terms.

I submit my academic CV and this cover letter through the NUS Careers portal.\\
You can reach me at \href{mailto:khsu@connect.ust.hk}{khsu@connect.ust.hk} or +852 6840 9753.

\vspace{0.2em}
Sincerely,

\vspace{0.2em}
Ke Cheng HSU \\
Ph.D.\ Candidate, Department of Economics, HKUST

\end{document}
